\documentclass[11pt, letterpaper]{article}
\input{/home/hyperion/.config/LatexHub/preambles/base.tex}
\input{/home/hyperion/.config/LatexHub/preambles/tikz.tex}
\input{/home/hyperion/.config/LatexHub/preambles/diagrams.tex}
\input{/home/hyperion/.config/LatexHub/preambles/macros-cat-theory.tex}
\input{/home/hyperion/.config/LatexHub/preambles/macros-algebra.tex}
\input{/home/hyperion/.config/LatexHub/preambles/basic-theorems-section.tex}
\input{/home/hyperion/.config/LatexHub/preambles/refs-basic.tex}
\input{/home/hyperion/.config/LatexHub/preambles/homework-formatting.tex}
\usepackage{titlepageBU}
\title{Homework 9}
\begin{document}
\makeproblem
\setcounter{section}{4}
\section{Problem 5}
\section{Problem 6}
\begin{problem}
	Let $R$ be a commutative ring with subset $S\subseteq R$ closed under multiplication and containing 1. Define the relation $\sim $ on $R\times S$ by, for all $s,s'\in S$ and $a,a'\in R$:
	\[
		(a,s)\sim (a',s') \iff \exists t\in S \text{ s.t. } t(s'a-sa')=0
	.\]
	Prove that $\sim $ is an equivalence relation.
\end{problem}
\begin{proof}
	For this problem, let $a,a',a''\in R$ and $s,s',s'',t,u\in S$. Given $(a,s)\sim (a',s')$, say the relation is \emph{witnessed by $t$} if
	\[
		t(as'-a's)=0
	.\]
	Of course the witness is not necessarily unique; it is simply convenient terminology.

	(reflexivity) By commutativity of $R$,
	\[
		1(as-sa)=as-as=0
	.\]
	Therefore, $(a,s)\sim (a,s)$.
	
	(symmetry) Let $(a,s)\sim (a',s')$ be witnessed by $t$. Then by commutativity of $R$,
	\[
		t(as'-a's)=0 \iff (-t)(a's-as')=0 \iff (a',s')\sim (a,s)
	,\]
	where the last implication follows by definition of $\sim $. Therefore, $\sim $ is symmetric.

	(transitivity) Let $(a,s)\sim (a',s')$ and $(a',s')\sim (a'',s'')$ be witnessed by $t$ and $u$, respectively. Then
	\[
		t(as'-a's) = 0 = u(a's''-a''s')
	.\]
	Since $u,s''\in S$, we have $us''\in S$, so we can write
	\[
		(us''t)(as'-a's)=0
	.\]
	a
\end{proof}
\begin{problem}
	Denote by $\frac{a}{s} $ the equivalence class $[(a,s)]_\sim $. Let $S^{-1} R =(R\times S)/\tildesim $ be endowed with the operations
	\[
		\frac{a}{s} +\frac{a'}{s'} =\frac{as'+a's}{ss'} ,\qquad \frac{a}{s} \cdot \frac{a'}{s'} =\frac{aa'}{ss'} 
	\]
	for all $a,a'\in R$ and $s,s'\in S$. Show that $S^{-1} R$ is a commutative ring.
\end{problem}
\begin{proof}
	It is immediate by definition that $S^{-1} R$ is closed under both operations. Since $R$ is commutative,
	\[
		\frac{a}{s} +\frac{a'}{s'} =\frac{as'+a's}{ss'} =\frac{a's+as'}{s's} =\frac{a'}{s'} +\frac{a}{s} 
	.\]
	We can also show associativity of addition:
	\[
		\left( \frac{a}{s} +\frac{a'}{s'}  \right) +\frac{a''}{s''} =\frac{as'+a's}{ss'} +\frac{a''}{s''} =\frac{\left( as'+a's \right) s''+a''ss'}{ss's''} =\frac{ass''+a'ss''+a''ss'}{ss's''} ,
	\]
	\[
		\frac{a}{s} +\left( \frac{a'}{s'} +\frac{a''}{s''}  \right) =\frac{a}{s} +\frac{a's''+a''s'}{s's''} =\frac{as's''+\left( a's''+a''s' \right) s}{ss's''} =\frac{as's''+a's''s+a''s's}{ss's''} 
	.\]
	Invoking commutativity of $R$ shows these quantities are equal. Since $S$ contains $1_R\in S$, define $1$ in $S^{-1} R$ as $1_R/1_R$. Then
	\[
		\frac{1_R}{1_R} \cdot \frac{a}{s} =\frac{1_Ra}{1_Rs} =\frac{a}{s} 
	\]
	since $1_R $ is the identity in $R$. To show commutativity, we note that since $R$ is commutative,
	\[
		\frac{a}{s} \cdot \frac{a'}{s'} =\frac{aa'}{ss'} =\frac{a'a}{s's} =\frac{a'}{s'} \cdot \frac{a}{s} 
	.\]
	The final thing we must do is deal with zero. Define $0$ as any element of the form $0/s$; we must show these all fall within the same equivalence class. Indeed, if $(0,s),(0,s')$ are two elements of $R\times S$, then
	\[
		(0s'-0s)=0-0=0,
	\]
	so we have $(0,s)\sim (0,s')$. Since the definition is now well-defined, we observe that
	\[
		\frac{a}{s} +\frac{0}{s'} = \frac{as'+0s}{ss'} =\frac{as'}{ss'} =\frac{a}{s} \cdot \frac{s'}{s'} 
	.\]
	It suffices to show that $(s',s')\sim (1_R,1_R)$. Once again,
	\[
		(s'1_R-1_Rs')=s'-s'=0
	,\]
	so we indeed have $s'/s'=1$. Therefore,
	\[
		\frac{a}{s} \cdot \frac{s'}{s'} =\frac{a}{s} \cdot 1=\frac{a}{s} 
	.\]
	0 is thus an additive identity. Given $a/s$, define the inverse $-a/s$ by $[(-a,s)]_{\tildesim}$. Then
	\[
		\frac{a}{s} +\frac{-a}{s} = \frac{-as + as}{s^2} =\frac{0}{s^2} =0
	.\]
	Therefore, $S^{-1} R$ is a commutative ring.
\end{proof}

\begin{problem}
	Show $S^{-1} R\cong 0$ if and only if $0\in S$.
\end{problem}
\begin{proof}
	It suffices to show that $(a,s)\sim (0,s')$ for all $(a,s)\in R\times S$ and some $s'\in S$. Consider $s'=0$. We have
	\[
		(a0-0s)=0-0=0,
	\]
	so indeed $(a,s)\sim (0,0)=0$. Therefore, $S^{-1} R$ is the zero ring.
\end{proof}


\section{Problem 7}



\end{document}
